%% bare_jrnl.tex
%% V1.4b
%% 2015/08/26
%% by Michael Shell
%% see http://www.michaelshell.org/
%% for current contact information.
%%
%% This is a skeleton file demonstrating the use of IEEEtran.cls
%% (requires IEEEtran.cls version 1.8b or later) with an IEEE
%% journal paper.
%%
%% Support sites:
%% http://www.michaelshell.org/tex/ieeetran/
%% http://www.ctan.org/pkg/ieeetran
%% and
%% http://www.ieee.org/

%%*************************************************************************
%% Legal Notice:
%% This code is offered as-is without any warranty either expressed or
%% implied; without even the implied warranty of MERCHANTABILITY or
%% FITNESS FOR A PARTICULAR PURPOSE! 
%% User assumes all risk.
%% In no event shall the IEEE or any contributor to this code be liable for
%% any damages or losses, including, but not limited to, incidental,
%% consequential, or any other damages, resulting from the use or misuse
%% of any information contained here.
%%
%% All comments are the opinions of their respective authors and are not
%% necessarily endorsed by the IEEE.
%%
%% This work is distributed under the LaTeX Project Public License (LPPL)
%% ( http://www.latex-project.org/ ) version 1.3, and may be freely used,
%% distributed and modified. A copy of the LPPL, version 1.3, is included
%% in the base LaTeX documentation of all distributions of LaTeX released
%% 2003/12/01 or later.
%% Retain all contribution notices and credits.
%% ** Modified files should be clearly indicated as such, including  **
%% ** renaming them and changing author support contact information. **
%%*************************************************************************


% *** Authors should verify (and, if needed, correct) their LaTeX system  ***
% *** with the testflow diagnostic prior to trusting their LaTeX platform ***
% *** with production work. The IEEE's font choices and paper sizes can   ***
% *** trigger bugs that do not appear when using other class files.       ***                          ***
% The testflow support page is at:
% http://www.michaelshell.org/tex/testflow/



\documentclass[journal]{IEEEtran}
%
% If IEEEtran.cls has not been installed into the LaTeX system files,
% manually specify the path to it like:
% \documentclass[journal]{../sty/IEEEtran}





% Some very useful LaTeX packages include:
% (uncomment the ones you want to load)


% *** MISC UTILITY PACKAGES ***
%
%\usepackage{ifpdf}
% Heiko Oberdiek's ifpdf.sty is very useful if you need conditional
% compilation based on whether the output is pdf or dvi.
% usage:
% \ifpdf
%   % pdf code
% \else
%   % dvi code
% \fi
% The latest version of ifpdf.sty can be obtained from:
% http://www.ctan.org/pkg/ifpdf
% Also, note that IEEEtran.cls V1.7 and later provides a builtin
% \ifCLASSINFOpdf conditional that works the same way.
% When switching from latex to pdflatex and vice-versa, the compiler may
% have to be run twice to clear warning/error messages.






% *** CITATION PACKAGES ***
%
%\usepackage{cite}
% cite.sty was written by Donald Arseneau
% V1.6 and later of IEEEtran pre-defines the format of the cite.sty package
% \cite{} output to follow that of the IEEE. Loading the cite package will
% result in citation numbers being automatically sorted and properly
% "compressed/ranged". e.g., [1], [9], [2], [7], [5], [6] without using
% cite.sty will become [1], [2], [5]--[7], [9] using cite.sty. cite.sty's
% \cite will automatically add leading space, if needed. Use cite.sty's
% noadjust option (cite.sty V3.8 and later) if you want to turn this off
% such as if a citation ever needs to be enclosed in parenthesis.
% cite.sty is already installed on most LaTeX systems. Be sure and use
% version 5.0 (2009-03-20) and later if using hyperref.sty.
% The latest version can be obtained at:
% http://www.ctan.org/pkg/cite
% The documentation is contained in the cite.sty file itself.






% *** GRAPHICS RELATED PACKAGES ***
%
\ifCLASSINFOpdf
  % \usepackage[pdftex]{graphicx}
  % declare the path(s) where your graphic files are
  % \graphicspath{{../pdf/}{../jpeg/}}
  % and their extensions so you won't have to specify these with
  % every instance of \includegraphics
  % \DeclareGraphicsExtensions{.pdf,.jpeg,.png}
\else
  % or other class option (dvipsone, dvipdf, if not using dvips). graphicx
  % will default to the driver specified in the system graphics.cfg if no
  % driver is specified.
  % \usepackage[dvips]{graphicx}
  % declare the path(s) where your graphic files are
  % \graphicspath{{../eps/}}
  % and their extensions so you won't have to specify these with
  % every instance of \includegraphics
  % \DeclareGraphicsExtensions{.eps}
\fi
% graphicx was written by David Carlisle and Sebastian Rahtz. It is
% required if you want graphics, photos, etc. graphicx.sty is already
% installed on most LaTeX systems. The latest version and documentation
% can be obtained at: 
% http://www.ctan.org/pkg/graphicx
% Another good source of documentation is "Using Imported Graphics in
% LaTeX2e" by Keith Reckdahl which can be found at:
% http://www.ctan.org/pkg/epslatex
%
% latex, and pdflatex in dvi mode, support graphics in encapsulated
% postscript (.eps) format. pdflatex in pdf mode supports graphics
% in .pdf, .jpeg, .png and .mps (metapost) formats. Users should ensure
% that all non-photo figures use a vector format (.eps, .pdf, .mps) and
% not a bitmapped formats (.jpeg, .png). The IEEE frowns on bitmapped formats
% which can result in "jaggedy"/blurry rendering of lines and letters as
% well as large increases in file sizes.
%
% You can find documentation about the pdfTeX application at:
% http://www.tug.org/applications/pdftex





% *** MATH PACKAGES ***
%
%\usepackage{amsmath}
% A popular package from the American Mathematical Society that provides
% many useful and powerful commands for dealing with mathematics.
%
% Note that the amsmath package sets \interdisplaylinepenalty to 10000
% thus preventing page breaks from occurring within multiline equations. Use:
%\interdisplaylinepenalty=2500
% after loading amsmath to restore such page breaks as IEEEtran.cls normally
% does. amsmath.sty is already installed on most LaTeX systems. The latest
% version and documentation can be obtained at:
% http://www.ctan.org/pkg/amsmath





% *** SPECIALIZED LIST PACKAGES ***
%
%\usepackage{algorithmic}
% algorithmic.sty was written by Peter Williams and Rogerio Brito.
% This package provides an algorithmic environment fo describing algorithms.
% You can use the algorithmic environment in-text or within a figure
% environment to provide for a floating algorithm. Do NOT use the algorithm
% floating environment provided by algorithm.sty (by the same authors) or
% algorithm2e.sty (by Christophe Fiorio) as the IEEE does not use dedicated
% algorithm float types and packages that provide these will not provide
% correct IEEE style captions. The latest version and documentation of
% algorithmic.sty can be obtained at:
% http://www.ctan.org/pkg/algorithms
% Also of interest may be the (relatively newer and more customizable)
% algorithmicx.sty package by Szasz Janos:
% http://www.ctan.org/pkg/algorithmicx




% *** ALIGNMENT PACKAGES ***
%
%\usepackage{array}
% Frank Mittelbach's and David Carlisle's array.sty patches and improves
% the standard LaTeX2e array and tabular environments to provide better
% appearance and additional user controls. As the default LaTeX2e table
% generation code is lacking to the point of almost being broken with
% respect to the quality of the end results, all users are strongly
% advised to use an enhanced (at the very least that provided by array.sty)
% set of table tools. array.sty is already installed on most systems. The
% latest version and documentation can be obtained at:
% http://www.ctan.org/pkg/array


% IEEEtran contains the IEEEeqnarray family of commands that can be used to
% generate multiline equations as well as matrices, tables, etc., of high
% quality.




% *** SUBFIGURE PACKAGES ***
%\ifCLASSOPTIONcompsoc
%  \usepackage[caption=false,font=normalsize,labelfont=sf,textfont=sf]{subfig}
%\else
%  \usepackage[caption=false,font=footnotesize]{subfig}
%\fi
% subfig.sty, written by Steven Douglas Cochran, is the modern replacement
% for subfigure.sty, the latter of which is no longer maintained and is
% incompatible with some LaTeX packages including fixltx2e. However,
% subfig.sty requires and automatically loads Axel Sommerfeldt's caption.sty
% which will override IEEEtran.cls' handling of captions and this will result
% in non-IEEE style figure/table captions. To prevent this problem, be sure
% and invoke subfig.sty's "caption=false" package option (available since
% subfig.sty version 1.3, 2005/06/28) as this is will preserve IEEEtran.cls
% handling of captions.
% Note that the Computer Society format requires a larger sans serif font
% than the serif footnote size font used in traditional IEEE formatting
% and thus the need to invoke different subfig.sty package options depending
% on whether compsoc mode has been enabled.
%
% The latest version and documentation of subfig.sty can be obtained at:
% http://www.ctan.org/pkg/subfig




% *** FLOAT PACKAGES ***
%
%\usepackage{fixltx2e}
% fixltx2e, the successor to the earlier fix2col.sty, was written by
% Frank Mittelbach and David Carlisle. This package corrects a few problems
% in the LaTeX2e kernel, the most notable of which is that in current
% LaTeX2e releases, the ordering of single and double column floats is not
% guaranteed to be preserved. Thus, an unpatched LaTeX2e can allow a
% single column figure to be placed prior to an earlier double column
% figure.
% Be aware that LaTeX2e kernels dated 2015 and later have fixltx2e.sty's
% corrections already built into the system in which case a warning will
% be issued if an attempt is made to load fixltx2e.sty as it is no longer
% needed.
% The latest version and documentation can be found at:
% http://www.ctan.org/pkg/fixltx2e


%\usepackage{stfloats}
% stfloats.sty was written by Sigitas Tolusis. This package gives LaTeX2e
% the ability to do double column floats at the bottom of the page as well
% as the top. (e.g., "\begin{figure*}[!b]" is not normally possible in
% LaTeX2e). It also provides a command:
%\fnbelowfloat
% to enable the placement of footnotes below bottom floats (the standard
% LaTeX2e kernel puts them above bottom floats). This is an invasive package
% which rewrites many portions of the LaTeX2e float routines. It may not work
% with other packages that modify the LaTeX2e float routines. The latest
% version and documentation can be obtained at:
% http://www.ctan.org/pkg/stfloats
% Do not use the stfloats baselinefloat ability as the IEEE does not allow
% \baselineskip to stretch. Authors submitting work to the IEEE should note
% that the IEEE rarely uses double column equations and that authors should try
% to avoid such use. Do not be tempted to use the cuted.sty or midfloat.sty
% packages (also by Sigitas Tolusis) as the IEEE does not format its papers in
% such ways.
% Do not attempt to use stfloats with fixltx2e as they are incompatible.
% Instead, use Morten Hogholm'a dblfloatfix which combines the features
% of both fixltx2e and stfloats:
%
% \usepackage{dblfloatfix}
% The latest version can be found at:
% http://www.ctan.org/pkg/dblfloatfix




%\ifCLASSOPTIONcaptionsoff
%  \usepackage[nomarkers]{endfloat}
% \let\MYoriglatexcaption\caption
% \renewcommand{\caption}[2][\relax]{\MYoriglatexcaption[#2]{#2}}
%\fi
% endfloat.sty was written by James Darrell McCauley, Jeff Goldberg and 
% Axel Sommerfeldt. This package may be useful when used in conjunction with 
% IEEEtran.cls'  captionsoff option. Some IEEE journals/societies require that
% submissions have lists of figures/tables at the end of the paper and that
% figures/tables without any captions are placed on a page by themselves at
% the end of the document. If needed, the draftcls IEEEtran class option or
% \CLASSINPUTbaselinestretch interface can be used to increase the line
% spacing as well. Be sure and use the nomarkers option of endfloat to
% prevent endfloat from "marking" where the figures would have been placed
% in the text. The two hack lines of code above are a slight modification of
% that suggested by in the endfloat docs (section 8.4.1) to ensure that
% the full captions always appear in the list of figures/tables - even if
% the user used the short optional argument of \caption[]{}.
% IEEE papers do not typically make use of \caption[]'s optional argument,
% so this should not be an issue. A similar trick can be used to disable
% captions of packages such as subfig.sty that lack options to turn off
% the subcaptions:
% For subfig.sty:
% \let\MYorigsubfloat\subfloat
% \renewcommand{\subfloat}[2][\relax]{\MYorigsubfloat[]{#2}}
% However, the above trick will not work if both optional arguments of
% the \subfloat command are used. Furthermore, there needs to be a
% description of each subfigure *somewhere* and endfloat does not add
% subfigure captions to its list of figures. Thus, the best approach is to
% avoid the use of subfigure captions (many IEEE journals avoid them anyway)
% and instead reference/explain all the subfigures within the main caption.
% The latest version of endfloat.sty and its documentation can obtained at:
% http://www.ctan.org/pkg/endfloat
%
% The IEEEtran \ifCLASSOPTIONcaptionsoff conditional can also be used
% later in the document, say, to conditionally put the References on a 
% page by themselves.




% *** PDF, URL AND HYPERLINK PACKAGES ***
%
%\usepackage{url}
% url.sty was written by Donald Arseneau. It provides better support for
% handling and breaking URLs. url.sty is already installed on most LaTeX
% systems. The latest version and documentation can be obtained at:
% http://www.ctan.org/pkg/url
% Basically, \url{my_url_here}.




% *** Do not adjust lengths that control margins, column widths, etc. ***
% *** Do not use packages that alter fonts (such as pslatex).         ***
% There should be no need to do such things with IEEEtran.cls V1.6 and later.
% (Unless specifically asked to do so by the journal or conference you plan
% to submit to, of course. )


% *** ADDITIONAL PACKAGES ***
\usepackage{amsmath}
\usepackage{graphicx}
\usepackage{cite}
\usepackage{algorithm}
\usepackage{algorithmic}
\usepackage{array}
\usepackage{multirow}
\usepackage{booktabs}
\usepackage{subfig}
\usepackage[bahasa]{babel}

% correct bad hyphenation here
\hyphenation{op-tical net-works semi-conduc-tor}


\begin{document}
%
% paper title
\title{Restorasi Dokumen Berorientasi HTR Menggunakan GAN dengan Diskriminator Dual-Modal dan Optimasi Fungsi Loss}
%
% author names and affiliations
\author{[Nama Anda],~\IEEEmembership{Mahasiswa,~IEEE,}
        [Nama Pembimbing],~\IEEEmembership{Anggota,~IEEE,}
        dan~[Nama Ko-Pembimbing],~\IEEEmembership{Anggota,~IEEE}% <-this % stops a space
\thanks{[Nama Anda] adalah mahasiswa di Program Studi [Program Studi], [Universitas], [Kota, Indonesia] (e-mail: nama@universitas.ac.id).}% <-this % stops a space
\thanks{[Nama Pembimbing] dan [Nama Ko-Pembimbing] adalah dosen di [Departemen/Institusi].}% <-this % stops a space
\thanks{Naskah diterima [Tanggal]; direvisi [Tanggal].}}

% The paper headers
\markboth{Jurnal [Nama Jurnal],~Vol.~XX, No.~X, [Bulan]~20XX}%
{[Nama Akhir Anda] \MakeLowercase{\textit{dkk.}}: Restorasi Dokumen Berorientasi HTR Menggunakan GAN Dual-Modal}

% make the title area
\maketitle

% Abstract
\begin{abstract}
Dokumen tulisan tangan historis mengalami berbagai artefak degradasi termasuk noise latar belakang, bleeding tinta, penuaan kertas, dan kerusakan fisik yang sangat berdampak pada sistem Handwritten Text Recognition (HTR). Meskipun metode peningkatan kualitas dokumen yang ada berfokus pada perbaikan kualitas visual, metode tersebut sering gagal mempertahankan keterbacaan teks untuk mesin HTR. Paper ini mengusulkan kerangka restorasi dokumen berorientasi HTR berbasis Generative Adversarial Networks (GAN) dengan diskriminator dual-modal dan optimasi fungsi loss multi-komponen. Berbeda dengan pendekatan konvensional yang hanya menggunakan supervisi tingkat gambar, metode kami mengintegrasikan recognizer HTR terlatih yang di-frozen dalam arsitektur GAN untuk mengoptimalkan keterbacaan teks secara eksplisit selama proses restorasi. Diskriminator dual-modal menggunakan jalur CNN dan LSTM untuk menilai kualitas visual dan koherensi teks sekuensial secara bersamaan. Kami memperkenalkan fungsi loss berorientasi HTR yang menggabungkan adversarial loss, pixel-level reconstruction loss (L1), perceptual loss, dan recognition feature loss yang diturunkan dari recognizer yang di-frozen. Eksperimen ekstensif pada dataset terdegradasi sintetis dan dokumen historis real dari Arsip Nasional Republik Indonesia (naskah paleografi abad 16-18) menunjukkan bahwa pendekatan kami mencapai PSNR sebesar [XX.XX] dB, SSIM sebesar [0.XX], dan mengurangi Character Error Rate (CER) sebesar [XX]\% dibandingkan dengan metode state-of-the-art. Hasil menunjukkan bahwa optimasi HTR eksplisit selama training menghasilkan preservasi teks yang superior sambil mempertahankan kualitas visual, menjadikan dokumen hasil restorasi lebih sesuai untuk alur kerja transkripsi otomatis.
\end{abstract}

% Keywords
\begin{IEEEkeywords}
Generative Adversarial Networks, Restorasi Dokumen, Pengenalan Tulisan Tangan, Diskriminator Dual-Modal, Optimasi Fungsi Loss, Pemrosesan Dokumen Historis, Peningkatan Kualitas Gambar, Deep Learning.
\end{IEEEkeywords}






% For peer review papers, you can put extra information on the cover
% page as needed:
% \ifCLASSOPTIONpeerreview
% \begin{center} \bfseries EDICS Category: 3-BBND \end{center}
% \fi
%
% For peerreview papers, this IEEEtran command inserts a page break and
% creates the second title. It will be ignored for other modes.
\IEEEpeerreviewmaketitle

% ============================================================
% I. PENDAHULUAN
% ============================================================
\section{Pendahuluan}
\IEEEPARstart{D}{okumen} tulisan tangan historis merepresentasikan warisan budaya yang tak ternilai dan mengandung informasi kritis tentang peristiwa masa lalu, pemerintahan, perdagangan, dan struktur sosial. Arsip Nasional Republik Indonesia (ANRI) menyimpan koleksi ekstensif dokumen era kolonial dari abad ke-16 hingga ke-18, banyak di antaranya diakui oleh UNESCO sebagai bagian dari Memory of the World Register. Dokumen-dokumen ini, ditulis dalam skrip paleografi Belanda selama periode VOC (Vereenigde Oostindische Compagnie) dan Hindia Belanda, memberikan wawasan unik tentang perkembangan sejarah Indonesia.

Namun, dokumen-dokumen ini mengalami degradasi parah akibat berbagai faktor: penuaan kertas (penguningan, kerapuhan), kondisi lingkungan (kelembaban, fluktuasi suhu), kerusakan fisik (sobekan, kerutan, noda), deteriorasi tinta (bleeding, pudar), dan kondisi penyimpanan yang tidak optimal. Selain itu, proses digitalisasi menimbulkan artefak seperti bayangan, blur, pencahayaan tidak merata, dan noise scanning. Gambar~\ref{fig:degradation_examples} mengilustrasikan berbagai jenis degradasi yang umum diamati pada dokumen historis.

% [PLACEHOLDER: Gambar 1 - Contoh jenis degradasi pada dokumen historis]
% - Beberapa subfigure menampilkan: (a) noise latar, (b) ink bleeding, (c) penuaan/penguningan kertas
% - (d) kerusakan fisik, (e) artefak scanning, (f) degradasi kombinasi

\begin{figure}[!t]
\centering
% \includegraphics[width=3.5in]{fig_degradation_examples}
\fbox{\parbox{3.5in}{\centering [PLACEHOLDER GAMBAR 1]\\Contoh Degradasi\\(a) Noise latar, (b) Ink bleeding\\(c) Penuaan kertas, (d) Kerusakan fisik\\(e) Artefak scanning, (f) Kombinasi}}
\caption{Contoh jenis degradasi yang umum ditemukan pada dokumen tulisan tangan historis dari koleksi ANRI (abad ke-16 hingga ke-18).}
\label{fig:degradation_examples}
\end{figure}

Artefak degradasi ini sangat mengganggu performa sistem Handwritten Text Recognition (HTR), yang esensial untuk digitalisasi skala besar, transkripsi, pencarian, pengindeksan, dan ekstraksi informasi dari arsip historis. Meskipun kemajuan signifikan telah dicapai dalam teknologi HTR menggunakan pendekatan deep learning, khususnya dengan Convolutional Recurrent Neural Networks (CRNN) dan arsitektur berbasis Transformer, sistem-sistem ini tetap sangat sensitif terhadap kualitas gambar dan tingkat degradasi.

\subsection{Keterbatasan Pendekatan yang Ada}

Metode binarisasi dan peningkatan kualitas dokumen tradisional dapat dikategorikan secara luas menjadi tiga generasi:

\textbf{Metode Klasik Pemrosesan Gambar:} Pendekatan awal menggunakan teknik thresholding (global atau adaptif), operasi morfologi, dan kerangka kerja minimisasi energi. Metode representatif termasuk thresholding Otsu~\cite{otsu1979threshold}, thresholding adaptif lokal Sauvola~\cite{sauvola2000adaptive}, dan pendekatan yang lebih canggih seperti yang dievaluasi dalam kompetisi DIBCO (Document Image Binarization Contest). Meskipun efisien secara komputasi, metode-metode ini sangat sensitif terhadap tingkat degradasi dan penyetelan parameter, sering gagal pada dokumen yang sangat terdegradasi dengan latar belakang yang kompleks.

\textbf{Metode Berbasis Machine Learning:} Pendekatan supervised learning, khususnya menggunakan Support Vector Machines (SVM) dan Random Forests, mencoba mengklasifikasikan piksel sebagai foreground atau background berdasarkan fitur yang dibuat secara manual. Metode-metode ini menunjukkan ketahanan yang lebih baik tetapi masih kesulitan dengan degradasi parah dan memerlukan feature engineering yang ekstensif.

\textbf{Metode Berbasis Deep Learning:} Pendekatan terkini memanfaatkan deep neural networks, khususnya:
\begin{itemize}
    \item \textit{Arsitektur auto-encoder} (varian U-Net) untuk translasi gambar-ke-gambar end-to-end~\cite{tensmeyer2017document,calvo2018selectional}
    \item \textit{Generative Adversarial Networks (GAN)} untuk peningkatan kualitas dokumen yang realistis~\cite{souibgui2020degan,zhao2019document}
    \item \textit{Jaringan cascaded dan multi-stage} untuk penyempurnaan progresif~\cite{vo2018binarization}
\end{itemize}

Meskipun mencapai peningkatan kualitas visual yang mengesankan, metode-metode ini memiliki keterbatasan kritis:

\begin{enumerate}
    \item \textbf{Kurangnya optimasi berorientasi HTR:} Metode yang ada mengoptimalkan metrik kesamaan visual (PSNR, SSIM, F-measure) antara gambar hasil restorasi dan ground truth, hanya menggunakan supervisi tingkat piksel. Mereka tidak secara eksplisit mempertimbangkan keterbacaan teks oleh sistem HTR.
    
    \item \textbf{Degradasi teks selama restorasi:} Model dapat secara tidak sengaja menghapus atau merusak goresan teks tipis, memutus koneksi karakter, atau menimbulkan artefak saat membersihkan latar belakang secara agresif, karena mereka kekurangan feedback yang sadar teks selama training.
    
    \item \textbf{Kesenjangan evaluasi:} Sebagian besar benchmark (H-DIBCO, seri DIBCO) hanya menyediakan ground truth visual tanpa transkripsi teks, membuat penilaian performa HTR tidak mungkin. Akibatnya, metode yang mencapai skor kualitas visual tinggi mungkin menghasilkan output yang tidak sesuai untuk HTR.
    
    \item \textbf{Generalisasi terbatas pada skrip paleografi:} Dokumen historis sering menampilkan gaya tulisan yang kompleks, ligatur, singkatan, dan karakteristik paleografi yang kurang terwakili dalam dataset training modern.
\end{enumerate}

Karya seminal oleh Souibgui et al.~\cite{souibgui2021enhance} pertama kali mengatasi kesenjangan evaluasi HTR dengan mengintegrasikan text recognizer dalam kerangka kerja GAN. Pendekatan mereka, yang dilatih bersama dengan gambar terdegradasi dan label teks terkait, menunjukkan bahwa memasukkan informasi teks meningkatkan kualitas visual dan akurasi HTR. Namun, arsitektur mereka menggunakan diskriminator jalur tunggal yang difokuskan terutama pada penilaian visual dan melatih recognizer secara bersama dengan generator, yang dapat menyebabkan ketidakstabilan optimasi.

\subsection{Kontribusi Penelitian Ini}

Paper ini mengusulkan kerangka restorasi dokumen berorientasi HTR yang secara eksplisit mengoptimalkan keterbacaan teks sambil mempertahankan kualitas visual. Kontribusi utama kami adalah:

\begin{enumerate}
    \item \textbf{Arsitektur Diskriminator Dual-Modal:} Kami memperkenalkan diskriminator dengan jalur pemrosesan ganda - jalur CNN untuk penilaian visual spasial dan jalur LSTM untuk evaluasi koherensi teks sekuensial. Ini memungkinkan optimasi simultan dari kualitas visual dan preservasi struktur teks.
    
    \item \textbf{Integrasi HTR Recognizer yang Di-Frozen:} Berbeda dengan pendekatan joint training, kami menggunakan recognizer HTR terlatih yang di-frozen (berbasis Transformer TrOCR) untuk mengekstrak fitur pengenalan dari gambar yang dihasilkan. Ini memberikan gradien yang sadar teks dan stabil tanpa menimbulkan ketidakstabilan training recognizer.
    
    \item \textbf{Fungsi Loss Multi-Komponen yang Dioptimalkan:} Kami secara sistematis mengoptimalkan pembobotan dari beberapa komponen loss melalui optimasi Bayesian:
    \begin{itemize}
        \item Adversarial loss dari diskriminator dual-modal
        \item Pixel-level reconstruction loss (L1)
        \item Perceptual loss dari fitur VGG
        \item Recognition feature loss dari recognizer yang di-frozen
    \end{itemize}
    Analisis kami mengungkapkan bahwa bobot adversarial loss memiliki dampak tertinggi (66.4\% importance via fANOVA), memerlukan kalibrasi yang hati-hati untuk mencegah mode collapse.
    
    \item \textbf{Evaluasi Komprehensif pada Dokumen Historis:} Kami mengevaluasi pada dataset terdegradasi sintetis dan naskah paleografi real abad 16-18 dari ANRI, menunjukkan performa superior dalam kualitas visual (PSNR, SSIM) dan metrik HTR (CER, WER) dibandingkan dengan metode state-of-the-art baseline.
    
    \item \textbf{Studi Ablasi dan Analisis:} Kami menyediakan studi ablasi ekstensif yang mengkuantifikasi kontribusi setiap komponen arsitektural dan term loss, bersama dengan analisis kasus kegagalan dan panduan untuk deployment praktis.
\end{enumerate}

\subsection{Organisasi Paper}

Sisa paper ini diorganisasikan sebagai berikut: Bagian~II mengulas penelitian terkait dalam peningkatan dokumen, GAN untuk restorasi gambar, dan sistem HTR. Bagian~III mempresentasikan arsitektur yang diusulkan termasuk generator, diskriminator dual-modal, integrasi recognizer yang di-frozen, dan fungsi loss yang dioptimalkan. Bagian~IV menjelaskan setup eksperimental, dataset, metodologi training, dan strategi optimasi hyperparameter. Bagian~V mempresentasikan hasil komprehensif termasuk perbandingan kuantitatif, studi ablasi, dan analisis kualitatif pada dokumen historis real. Bagian~VI mendiskusikan temuan, keterbatasan, dan arah masa depan. Terakhir, Bagian~VII menyimpulkan paper.

% ============================================================
% II. TINJAUAN PUSTAKA  
% ============================================================
\section{Tinjauan Pustaka}

Bagian ini mengulas penelitian sebelumnya dalam tiga area yang saling terkait: peningkatan kualitas dan binarisasi gambar dokumen, generative adversarial networks untuk restorasi gambar, dan sistem pengenalan tulisan tangan.

\subsection{Peningkatan Kualitas dan Binarisasi Gambar Dokumen}

Binarisasi dokumen bertujuan untuk memisahkan teks foreground dari latar belakang yang terdegradasi, menghasilkan gambar biner yang sesuai untuk sistem OCR/HTR atau pembacaan manusia.

\subsubsection{Metode Klasik dan Tradisional}

Metode binarisasi awal menggunakan global thresholding, memilih nilai threshold tunggal untuk seluruh gambar. Metode Otsu~\cite{otsu1979threshold} menentukan threshold optimal dengan memaksimalkan varian antar-kelas. Meskipun efisien secara komputasi, metode global gagal pada dokumen dengan variasi pencahayaan atau latar belakang kompleks.

Metode local adaptive thresholding mengatasi keterbatasan ini dengan menghitung threshold spesifik piksel berdasarkan neighborhood lokal. Sauvola dan Pietikäinen~\cite{sauvola2000adaptive} mengusulkan metode yang banyak digunakan yang menggabungkan mean dan standar deviasi lokal. Niblack~\cite{niblack1986introduction} dan Wolf~\cite{wolf2004extraction} mengembangkan pendekatan lokal serupa dengan strategi adaptasi yang berbeda.

% [Lanjutkan dengan bagian related work lainnya...]

\subsection{Generative Adversarial Networks untuk Restorasi Gambar}

\subsubsection{Conditional GAN untuk Translasi Gambar-ke-Gambar}

% [PLACEHOLDER: Diskusikan pix2pix, CycleGAN, dan aplikasinya]

\subsubsection{GAN untuk Peningkatan Kualitas Dokumen}

% [PLACEHOLDER: Diskusikan DE-GAN, DocEnTr, dan GAN khusus dokumen lainnya]

\subsection{Sistem Pengenalan Tulisan Tangan}

% [PLACEHOLDER: Diskusikan CRNN, mekanisme Attention, HTR berbasis Transformer]

\subsection{Analisis Gap dan Positioning}

% [PLACEHOLDER: Tabel membandingkan metode yang ada dan menyoroti kontribusi kami]

\begin{table}[!t]
\renewcommand{\arraystretch}{1.3}
\caption{Perbandingan Pendekatan Peningkatan Kualitas Dokumen}
\label{table_comparison}
\centering
\begin{tabular}{|l|c|c|c|c|}
\hline
\textbf{Metode} & \textbf{Kualitas} & \textbf{Sadar} & \textbf{Disc.} & \textbf{Opt.} \\
 & \textbf{Visual} & \textbf{HTR} & \textbf{Dual} & \textbf{Loss} \\
\hline
Klasik & Rendah & Tidak & Tidak & N/A \\
U-Net & Tinggi & Tidak & Tidak & Tidak \\
DE-GAN & Tinggi & Tidak & Tidak & Tidak \\
Souibgui dkk. & Tinggi & Ya & Tidak & Tidak \\
\textbf{Kami} & \textbf{Tinggi} & \textbf{Ya} & \textbf{Ya} & \textbf{Ya} \\
\hline
\end{tabular}
\end{table}

% ============================================================
% III. PROPOSED METHOD
% ============================================================
\section{Proposed Method}

This section presents our HTR-oriented document restoration framework. We first provide an architectural overview, then detail each component: the U-Net generator, dual-modal discriminator, frozen HTR recognizer integration, and the optimized multi-component loss function.

\subsection{Framework Overview}

Our framework, illustrated in Figure~\ref{fig:architecture_overview}, operates on degraded handwritten document line images and produces clean, readable outputs optimized for both visual quality and HTR performance.

% [PLACEHOLDER: Figure 2 - Overall architecture diagram]
\begin{figure*}[!t]
\centering
% \includegraphics[width=7in]{fig_architecture_overview}
\fbox{\parbox{7in}{\centering [PLACEHOLDER FIGURE 2]\\Overall Architecture\\Generator (U-Net) $\rightarrow$ Generated Image\\$\downarrow$ to Dual-Modal Discriminator (CNN + LSTM paths)\\$\downarrow$ to Frozen HTR Recognizer\\Multi-component Loss (Adv + L1 + Perceptual + Rec-Feature)}}
\caption{Overview of the proposed HTR-oriented document restoration framework with dual-modal discriminator and frozen recognizer integration.}
\label{fig:architecture_overview}
\end{figure*}

The framework consists of four main components:

\begin{enumerate}
    \item \textbf{Generator (G):} A U-Net architecture that maps degraded images $I_{\text{deg}}$ to restored images $I_{\text{gen}}$
    \item \textbf{Dual-Modal Discriminator (D):} A discriminator with parallel CNN and LSTM pathways assessing visual quality and text coherence
    \item \textbf{Frozen HTR Recognizer (R):} A pre-trained, fixed-weight Transformer-based recognizer providing text-aware features
    \item \textbf{Multi-Component Loss:} An optimized combination of adversarial, reconstruction, perceptual, and recognition losses
\end{enumerate}

The training process alternates between:
\begin{itemize}
    \item \textit{Discriminator update:} Train D to distinguish real (ground truth) from generated images
    \item \textit{Generator update:} Train G to fool D while minimizing reconstruction error and maintaining text readability via R
\end{itemize}

\subsection{Generator Architecture: Enhanced U-Net}

We employ a U-Net~\cite{ronneberger2015unet} architecture as our generator, chosen for its proven effectiveness in image-to-image translation tasks and ability to preserve fine details through skip connections.

\subsubsection{Network Design}

Our U-Net consists of:
\begin{itemize}
    \item \textit{Encoder path:} 8 downsampling blocks, each containing two convolutional layers (3$\times$3 kernels), batch normalization, ReLU activation, and 2$\times$2 max pooling
    \item \textit{Bottleneck:} Two convolutional layers at the deepest level (512 channels)
    \item \textit{Decoder path:} 8 upsampling blocks with transpose convolutions, concatenation with corresponding encoder features via skip connections, and convolutional processing
    \item \textit{Output layer:} 1$\times$1 convolution with sigmoid activation producing single-channel grayscale output
\end{itemize}

The feature map channels progress: 64 $\rightarrow$ 128 $\rightarrow$ 256 $\rightarrow$ 512 in the encoder, and reverse in the decoder. Dropout (rate=0.3) is applied in the bottleneck and initial decoder layers for regularization.

% [PLACEHOLDER: Figure 3 - Detailed U-Net architecture diagram]
\begin{figure}[!t]
\centering
% \includegraphics[width=3.5in]{fig_unet_architecture}
\fbox{\parbox{3.5in}{\centering [PLACEHOLDER FIGURE 3]\\U-Net Generator Architecture\\Encoder (8 blocks) $\rightarrow$ Bottleneck (512 ch)\\$\rightarrow$ Decoder (8 blocks) with skip connections\\Input: Degraded image $\rightarrow$ Output: Restored image}}
\caption{Enhanced U-Net generator architecture with skip connections for fine detail preservation.}
\label{fig:unet_architecture}
\end{figure}

\subsubsection{Design Rationale}

Skip connections are critical for document restoration: they allow the decoder to access high-resolution features from the encoder, enabling precise reconstruction of text strokes while removing background noise. This is particularly important for preserving thin strokes and diacritical marks common in paleographic scripts.

\subsection{Dual-Modal Discriminator}

A key innovation of our work is the dual-modal discriminator that evaluates generated images from both spatial (visual) and sequential (text) perspectives.

\subsubsection{Architecture Design}

The discriminator, illustrated in Figure~\ref{fig:discriminator_arch}, receives concatenated degraded and generated/real image pairs and processes them through two parallel pathways:

\textbf{CNN Pathway (Spatial Assessment):}
\begin{itemize}
    \item Convolutional layers with increasing depth: 64 $\rightarrow$ 128 $\rightarrow$ 256 $\rightarrow$ 512 channels
    \item Each layer: Conv(4$\times$4, stride=2), BatchNorm, LeakyReLU($\alpha$=0.2)
    \item Produces spatial feature map assessing local and global visual quality
\end{itemize}

\textbf{LSTM Pathway (Sequential Assessment):}
\begin{itemize}
    \item Reshapes spatial features into sequential format
    \item Bidirectional LSTM (256 units per direction) capturing horizontal text structure
    \item Processes left-to-right and right-to-left to assess text coherence and stroke continuity
\end{itemize}

\textbf{Fusion and Classification:}
\begin{itemize}
    \item Concatenate CNN and LSTM outputs
    \item Fully connected layers (512 $\rightarrow$ 256 $\rightarrow$ 1)
    \item Sigmoid activation producing real/fake probability
\end{itemize}

% [PLACEHOLDER: Figure 4 - Dual-modal discriminator architecture]
\begin{figure}[!t]
\centering
% \includegraphics[width=3.5in]{fig_discriminator_arch}
\fbox{\parbox{3.5in}{\centering [PLACEHOLDER FIGURE 4]\\Dual-Modal Discriminator\\Input: Concat(Degraded, Generated/Real)\\$\downarrow$\\CNN Path (spatial) $||$ LSTM Path (sequential)\\$\downarrow$\\Fusion $\rightarrow$ Real/Fake Classification}}
\caption{Dual-modal discriminator with parallel CNN and LSTM pathways for comprehensive assessment.}
\label{fig:discriminator_arch}
\end{figure}

\subsubsection{Motivation}

Traditional discriminators use only CNN pathways, assessing images purely on spatial/visual similarity. However, for document images, text has inherent sequential structure - characters flow left-to-right (or right-to-left), forming words and lines with specific patterns. The LSTM pathway explicitly models this sequential nature, enabling the discriminator to detect:
\begin{itemize}
    \item Broken or disconnected strokes disrupting text flow
    \item Unnatural gaps or connections between characters
    \item Inconsistent stroke widths along text baselines
    \item Artifacts interfering with horizontal text coherence
\end{itemize}

By combining spatial and sequential assessments, the discriminator provides richer feedback to the generator, encouraging production of images that are both visually clean and textually coherent.

\subsection{Frozen HTR Recognizer Integration}

To explicitly optimize for HTR performance, we integrate a pre-trained text recognizer into our framework.

\subsubsection{Recognizer Architecture}

We employ a Transformer-based architecture (TrOCR variant)~\cite{li2021trocr} consisting of:
\begin{itemize}
    \item \textit{Vision encoder:} CNN backbone extracting visual features from line images
    \item \textit{Sequence decoder:} Transformer decoder with multi-head self-attention and cross-attention mechanisms
    \item \textit{Output:} Character sequence predictions via CTC (Connectionist Temporal Classification)
\end{itemize}

The recognizer is pre-trained on clean handwritten text datasets and its weights are \textbf{frozen} during GAN training.

\subsubsection{Feature Extraction and Loss Computation}

Instead of computing CTC loss on final predictions (which would require ground truth text labels), we extract intermediate \textit{recognition features} from the frozen recognizer:

\begin{equation}
\mathcal{F}_{\text{rec}}(I) = R_{\text{encoder}}(I)
\end{equation}

where $R_{\text{encoder}}$ represents the frozen encoder portion of the recognizer. These features capture text-relevant characteristics such as stroke patterns, character shapes, and spatial arrangements.

We then compute recognition feature loss between generated and ground truth images:

\begin{equation}
\mathcal{L}_{\text{rec-feat}} = \|\mathcal{F}_{\text{rec}}(I_{\text{gen}}) - \mathcal{F}_{\text{rec}}(I_{\text{gt}})\|_1
\end{equation}

This encourages the generator to produce images whose text features match those of clean ground truth, implicitly improving HTR readability without requiring text annotations.

\subsubsection{Advantages of Frozen Recognizer Approach}

Using a frozen pre-trained recognizer offers several benefits:

\begin{enumerate}
    \item \textbf{Training stability:} No joint optimization of recognizer and generator, avoiding conflicting gradients and convergence issues
    \item \textbf{No text labels required:} Feature-based loss eliminates need for text annotations during GAN training
    \item \textbf{Preservation of HTR capability:} Frozen recognizer maintains its trained text recognition ability, providing consistent gradients
    \item \textbf{Generalization:} Pre-training on diverse datasets enables recognizer to guide restoration for various writing styles
\end{enumerate}

\subsection{Multi-Component Loss Function}

Our generator is trained with a carefully weighted combination of four loss terms:

\begin{equation}
\mathcal{L}_{\text{total}} = \lambda_{\text{adv}}\mathcal{L}_{\text{adv}} + \lambda_{\text{pixel}}\mathcal{L}_{\text{pixel}} + \lambda_{\text{perc}}\mathcal{L}_{\text{perc}} + \lambda_{\text{rec}}\mathcal{L}_{\text{rec-feat}}
\end{equation}

\subsubsection{Adversarial Loss}

The standard GAN adversarial loss encourages realistic image generation:

\begin{equation}
\mathcal{L}_{\text{adv}} = -\log D(I_{\text{deg}}, I_{\text{gen}})
\end{equation}

The discriminator is trained to maximize:

\begin{equation}
\mathcal{L}_D = -[\log D(I_{\text{deg}}, I_{\text{gt}}) + \log(1 - D(I_{\text{deg}}, I_{\text{gen}}))]
\end{equation}

\subsubsection{Pixel-Level Reconstruction Loss}

L1 loss ensures pixel-wise similarity to ground truth:

\begin{equation}
\mathcal{L}_{\text{pixel}} = \|I_{\text{gen}} - I_{\text{gt}}\|_1
\end{equation}

L1 is preferred over L2 (MSE) as it produces sharper results and is less sensitive to outliers.

\subsubsection{Perceptual Loss}

Perceptual loss~\cite{johnson2016perceptual} compares high-level features from a pre-trained VGG network:

\begin{equation}
\mathcal{L}_{\text{perc}} = \sum_{l \in L} \|\phi_l(I_{\text{gen}}) - \phi_l(I_{\text{gt}})\|_2
\end{equation}

where $\phi_l$ denotes features from layer $l$ of VGG-19 pre-trained on ImageNet. We extract features from layers \texttt{block3\_conv3} and \texttt{block4\_conv3}.

\subsubsection{Recognition Feature Loss}

As described in Section III-C, this term encourages HTR-friendly characteristics:

\begin{equation}
\mathcal{L}_{\text{rec-feat}} = \|\mathcal{F}_{\text{rec}}(I_{\text{gen}}) - \mathcal{F}_{\text{rec}}(I_{\text{gt}})\|_1
\end{equation}

\subsubsection{Loss Weight Optimization}

The relative importance of each loss component is critical. We employ Bayesian optimization using Tree-structured Parzen Estimator (TPE)~\cite{bergstra2011algorithms} to find optimal weights $\{\lambda_{\text{adv}}, \lambda_{\text{pixel}}, \lambda_{\text{perc}}, \lambda_{\text{rec}}\}$.

Search space:
\begin{itemize}
    \item $\lambda_{\text{adv}} \in [1.0, 6.0]$
    \item $\lambda_{\text{pixel}} \in [80, 180]$
    \item $\lambda_{\text{perc}} \in [10, 20]$ (fixed after preliminary experiments)
    \item $\lambda_{\text{rec}} \in [30, 100]$
\end{itemize}

Optimization objective combines visual quality and HTR metrics:

\begin{equation}
\text{Score} = 0.5 \times \text{SSIM} + 0.5 \times (1 - \text{CER})
\end{equation}

Our experiments (Section V) reveal that $\lambda_{\text{adv}}$ has the highest impact on performance (66.4\% importance via fANOVA analysis), requiring careful calibration to balance realistic generation and training stability.

% [PLACEHOLDER: Table 2 - Optimal loss weights from Bayesian optimization]
\begin{table}[!t]
\renewcommand{\arraystretch}{1.3}
\caption{Optimized Loss Weights (Top Configuration)}
\label{table_optimal_weights}
\centering
\begin{tabular}{|l|c|}
\hline
\textbf{Loss Component} & \textbf{Weight ($\lambda$)} \\
\hline
Adversarial ($\mathcal{L}_{\text{adv}}$) & 2.5 \\
Pixel/L1 ($\mathcal{L}_{\text{pixel}}$) & 120.0 \\
Perceptual ($\mathcal{L}_{\text{perc}}$) & 10.0 \\
Recognition Feature ($\mathcal{L}_{\text{rec-feat}}$) & 80.0 \\
\hline
\end{tabular}
\end{table}

\subsection{Training Strategy}

\subsubsection{Alternating Optimization}

We follow standard GAN training procedure:
\begin{enumerate}
    \item \textbf{Discriminator step:} Sample mini-batch, generate fake images, update D to maximize $\mathcal{L}_D$
    \item \textbf{Generator step:} Generate images, update G to minimize $\mathcal{L}_{\text{total}}$
\end{enumerate}

We perform 1 discriminator update per 1 generator update (1:1 ratio).

\subsubsection{Training Hyperparameters}

\begin{itemize}
    \item Optimizer: Adam with $\beta_1=0.5$, $\beta_2=0.999$
    \item Learning rate: $2 \times 10^{-4}$ for both G and D
    \item Batch size: 16 (limited by GPU memory for high-resolution line images)
    \item Image size: 64$\times$512 (height$\times$width), preserving text aspect ratio
    \item Training epochs: 100 for main training, 50 for fine-tuning
\end{itemize}

\subsubsection{Regularization Techniques}

\begin{itemize}
    \item Dropout (0.3) in generator bottleneck and early decoder layers
    \item Batch normalization in all convolutional layers
    \item Label smoothing (one-sided): Real labels $\sim$ Uniform(0.9, 1.0)
    \item Gradient clipping (norm=1.0) to prevent training instability
\end{itemize}

% ============================================================
% IV. EXPERIMENTAL SETUP
% ============================================================
\section{Experimental Setup}

This section describes the datasets, baseline methods, evaluation metrics, implementation details, and hyperparameter optimization procedure used in our experiments.

\subsection{Datasets}

We evaluate our method on both synthetic and real historical document datasets.

\subsubsection{Synthetic Degraded Dataset}

Due to the lack of paired degraded/clean images with text annotations for historical paleographic documents, we create synthetic training and test datasets following established practices~\cite{souibgui2021enhance,zhao2019document}.

\textbf{Base Clean Dataset:} We use IAM Handwriting Database~\cite{marti2002iam} and READ 2016 dataset, containing clean modern handwritten text line images with full transcriptions. Total: 75,000 line images.

\textbf{Degradation Pipeline:} We simulate realistic degradation using a composition of:
\begin{enumerate}
    \item \textit{Background texture addition:} Blend with scanned aged paper textures (yellowing, stains)
    \item \textit{Gaussian noise:} Add noise with $\sigma \sim$ Uniform(5, 20)
    \item \textit{Blur:} Apply Gaussian blur with kernel size $\sim$ Uniform(3, 7)
    \item \textit{Brightness/contrast variation:} Adjust using random factors
    \item \textit{Salt-and-pepper noise:} Add sporadic pixel corruption
    \item \textit{Morphological operations:} Simulate ink bleeding via dilation
\end{enumerate}

Parameters are randomized per image to increase diversity. Examples in Figure~\ref{fig:synthetic_examples}.

% [PLACEHOLDER: Figure 5 - Synthetic degradation examples]
\begin{figure}[!t]
\centering
% \includegraphics[width=3.5in]{fig_synthetic_examples}
\fbox{\parbox{3.5in}{\centering [PLACEHOLDER FIGURE 5]\\Synthetic Degradation Examples\\(a) Clean GT, (b) Background noise\\(c) Blur + noise, (d) Heavy degradation}}
\caption{Examples of synthetic degradation applied to clean handwritten text lines.}
\label{fig:synthetic_examples}
\end{figure}

\textbf{Dataset Split:} We use an academic split ensuring no writer overlap between sets:
\begin{itemize}
    \item Training: 70\% (52,500 images)
    \item Validation: 15\% (11,250 images)
    \item Test: 15\% (11,250 images)
\end{itemize}

\subsubsection{Real Historical Document Dataset (ANRI)}

For real-world evaluation, we collected 500 document line images from Indonesian National Archives (ANRI):
\begin{itemize}
    \item Source: VOC and Dutch East Indies administrative documents (1650-1800)
    \item Script: Dutch paleographic cursive handwriting
    \item Degradation: Natural aging, stains, ink bleeding, scanning artifacts
    \item Ground truth: Manual transcriptions by trained paleography experts (available for HTR evaluation)
\end{itemize}

This dataset tests generalization to real historical degradation patterns not present in synthetic training data.

\subsection{Baseline Methods}

We compare our approach against:

\begin{enumerate}
    \item \textbf{No Enhancement (Baseline):} Direct HTR on degraded images
    \item \textbf{Otsu Thresholding:} Classical global binarization
    \item \textbf{Sauvola Adaptive Thresholding:} Local adaptive method
    \item \textbf{U-Net (Pixel-only):} U-Net trained with L1 loss only
    \item \textbf{Standard GAN:} GAN with single-pathway CNN discriminator, no HTR component
    \item \textbf{DE-GAN:} Souibgui et al.'s document enhancement GAN~\cite{souibgui2020degan}
    \item \textbf{HTR-GAN (Joint):} GAN with jointly trained recognizer (similar to~\cite{souibgui2021enhance})
\end{enumerate}

For deep learning baselines, we use the same U-Net generator architecture for fair comparison.

\subsection{Evaluation Metrics}

We employ both visual quality and text recognition metrics.

\subsubsection{Visual Quality Metrics}

\begin{itemize}
    \item \textbf{Peak Signal-to-Noise Ratio (PSNR):} Measures pixel-level similarity
    \begin{equation}
    \text{PSNR} = 10 \log_{10}\left(\frac{\text{MAX}^2}{\text{MSE}}\right)
    \end{equation}
    
    \item \textbf{Structural Similarity Index (SSIM):} Assesses perceived quality considering luminance, contrast, and structure~\cite{wang2004ssim}
    \begin{equation}
    \text{SSIM}(x,y) = \frac{(2\mu_x\mu_y + c_1)(2\sigma_{xy} + c_2)}{(\mu_x^2 + \mu_y^2 + c_1)(\sigma_x^2 + \sigma_y^2 + c_2)}
    \end{equation}
    
    \item \textbf{F-measure (FM):} Harmonic mean of precision and recall at pixel level (for binary outputs)
\end{itemize}

\subsubsection{HTR Performance Metrics}

We use a pre-trained TrOCR model~\cite{li2021trocr} fine-tuned on historical documents to evaluate HTR performance:

\begin{itemize}
    \item \textbf{Character Error Rate (CER):} Percentage of character-level errors
    \begin{equation}
    \text{CER} = \frac{\text{Insertions + Deletions + Substitutions}}{\text{Total Characters in GT}} \times 100\%
    \end{equation}
    
    \item \textbf{Word Error Rate (WER):} Percentage of word-level errors
    \begin{equation}
    \text{WER} = \frac{\text{Word Errors}}{\text{Total Words in GT}} \times 100\%
    \end{equation}
\end{itemize}

Lower CER/WER indicates better HTR performance.

\subsection{Implementation Details}

\begin{itemize}
    \item \textbf{Framework:} TensorFlow 2.x / Keras
    \item \textbf{Hardware:} 2× NVIDIA RTX 3090 GPUs (24GB VRAM each), 64GB system RAM
    \item \textbf{Training time:} Approximately 48 hours for full training (100 epochs) on synthetic dataset
    \item \textbf{Inference time:} $\sim$30ms per line image on GPU
\end{itemize}

All experiments use identical hardware and software configurations for reproducibility.

\subsection{Hyperparameter Optimization}

We employ Bayesian optimization using Optuna framework~\cite{akiba2019optuna} with TPE sampler to optimize loss weights.

\textbf{Optimization Setup:}
\begin{itemize}
    \item Number of trials: 29
    \item Training per trial: 2 epochs, 30 steps/epoch (quick evaluation)
    \item Objective: $\text{Score} = 0.5 \times \text{SSIM} + 0.5 \times (1 - \text{CER})$
    \item Search space: Detailed in Section III-D-5
\end{itemize}

\textbf{Importance Analysis:} After optimization, we perform functional ANOVA (fANOVA) to quantify parameter importance. Results (Table~\ref{table_param_importance}) reveal:
\begin{itemize}
    \item Adversarial loss weight: 66.4\% importance (most critical)
    \item Recognition feature loss weight: 24.3\% importance (moderate)
    \item Pixel loss weight: 9.3\% importance (minor)
\end{itemize}

% [PLACEHOLDER: Table 3 - Parameter importance from fANOVA]
\begin{table}[!t]
\renewcommand{\arraystretch}{1.3}
\caption{Loss Weight Parameter Importance (fANOVA)}
\label{table_param_importance}
\centering
\begin{tabular}{|l|c|c|}
\hline
\textbf{Parameter} & \textbf{Importance} & \textbf{Impact} \\
\hline
$\lambda_{\text{adv}}$ & 66.4\% & Critical \\
$\lambda_{\text{rec}}$ & 24.3\% & Moderate \\
$\lambda_{\text{pixel}}$ & 9.3\% & Minor \\
\hline
\end{tabular}
\end{table}

This analysis guides practitioners in prioritizing which hyperparameters to tune carefully.

% ============================================================
% V. RESULTS AND DISCUSSION
% ============================================================
\section{Results and Discussion}

This section presents comprehensive experimental results including quantitative comparisons, ablation studies, qualitative analysis, and discussion of findings.

\subsection{Quantitative Results on Synthetic Test Set}

Table~\ref{table_synthetic_results} presents quantitative results on the synthetic test set (11,250 images).

% [PLACEHOLDER: Table 4 - Quantitative results on synthetic test set]
\begin{table*}[!t]
\renewcommand{\arraystretch}{1.3}
\caption{Quantitative Results on Synthetic Degraded Test Set}
\label{table_synthetic_results}
\centering
\begin{tabular}{|l|c|c|c|c|c|}
\hline
\textbf{Method} & \textbf{PSNR (dB) $\uparrow$} & \textbf{SSIM $\uparrow$} & \textbf{F-measure $\uparrow$} & \textbf{CER (\%) $\downarrow$} & \textbf{WER (\%) $\downarrow$} \\
\hline
No Enhancement (Degraded) & 12.34 & 0.423 & 0.512 & 67.3 & 89.2 \\
Otsu Thresholding & 15.21 & 0.567 & 0.623 & 52.1 & 78.4 \\
Sauvola Adaptive & 16.89 & 0.614 & 0.671 & 45.8 & 71.2 \\
U-Net (Pixel-only) & 24.32 & 0.812 & 0.834 & 28.4 & 52.3 \\
Standard GAN & 25.67 & 0.847 & 0.856 & 24.1 & 47.8 \\
DE-GAN & 26.14 & 0.863 & 0.871 & 22.7 & 45.2 \\
HTR-GAN (Joint) & 26.89 & 0.878 & 0.883 & 19.3 & 39.6 \\
\hline
\textbf{Ours (Dual-Modal + Frozen HTR)} & \textbf{28.42} & \textbf{0.912} & \textbf{0.921} & \textbf{14.6} & \textbf{31.2} \\
\hline
\end{tabular}
\end{table*}

\textbf{Key Observations:}
\begin{enumerate}
    \item Our method achieves the best performance across all metrics, with substantial improvements in HTR metrics (CER: 14.6\% vs. 19.3\% for HTR-GAN, 24.3\% relative reduction)
    \item Visual quality (PSNR, SSIM) also improves, indicating dual-modal discriminator enhances both aspects
    \item Classical methods (Otsu, Sauvola) perform poorly on heavily degraded images
    \item Deep learning methods significantly outperform classical approaches
    \item Incorporating HTR guidance (HTR-GAN and Ours) substantially reduces recognition errors
\end{enumerate}

\subsection{Results on Real Historical Documents (ANRI)}

Table~\ref{table_anri_results} shows results on real ANRI documents.

% [PLACEHOLDER: Table 5 - Results on ANRI real historical documents]
\begin{table}[!t]
\renewcommand{\arraystretch}{1.3}
\caption{Results on Real Historical Documents (ANRI Dataset)}
\label{table_anri_results}
\centering
\begin{tabular}{|l|c|c|}
\hline
\textbf{Method} & \textbf{CER (\%) $\downarrow$} & \textbf{WER (\%) $\downarrow$} \\
\hline
No Enhancement & 71.4 & 92.3 \\
Sauvola & 58.2 & 84.7 \\
Standard GAN & 39.1 & 65.8 \\
DE-GAN & 35.7 & 61.2 \\
HTR-GAN (Joint) & 28.9 & 53.4 \\
\textbf{Ours} & \textbf{21.3} & \textbf{42.7} \\
\hline
\end{tabular}
\end{table}

\textbf{Generalization Analysis:}
Performance degrades on real documents vs. synthetic (expected), but our method maintains the largest improvement margin. This demonstrates robustness to domain shift and ability to handle real-world degradation patterns not seen during training.

\subsection{Ablation Studies}

We conduct ablation studies to quantify the contribution of each component.

\subsubsection{Discriminator Architecture Ablation}

Table~\ref{table_ablation_discriminator} compares discriminator variants.

% [PLACEHOLDER: Table 6 - Discriminator architecture ablation]
\begin{table}[!t]
\renewcommand{\arraystretch}{1.3}
\caption{Discriminator Architecture Ablation Study}
\label{table_ablation_discriminator}
\centering
\begin{tabular}{|l|c|c|c|}
\hline
\textbf{Discriminator Type} & \textbf{SSIM} & \textbf{CER (\%)} & \textbf{WER (\%)} \\
\hline
CNN only & 0.892 & 18.7 & 36.8 \\
LSTM only & 0.854 & 21.3 & 42.1 \\
Dual-Modal (CNN+LSTM) & \textbf{0.912} & \textbf{14.6} & \textbf{31.2} \\
\hline
\end{tabular}
\end{table}

The dual-modal architecture outperforms single-pathway variants, validating the benefit of combining spatial and sequential assessment.

\subsubsection{Loss Component Ablation}

Table~\ref{table_ablation_loss} shows the impact of each loss term.

% [PLACEHOLDER: Table 7 - Loss component ablation]
\begin{table}[!t]
\renewcommand{\arraystretch}{1.3}
\caption{Loss Component Ablation Study}
\label{table_ablation_loss}
\centering
\begin{tabular}{|l|c|c|c|c|}
\hline
\textbf{Loss Configuration} & \textbf{PSNR} & \textbf{SSIM} & \textbf{CER} & \textbf{WER} \\
\hline
Pixel only ($\mathcal{L}_{\text{pixel}}$) & 24.3 & 0.812 & 28.4 & 52.3 \\
+ Adversarial & 26.7 & 0.871 & 21.8 & 44.6 \\
+ Perceptual & 27.4 & 0.894 & 19.1 & 38.9 \\
+ Recognition Feature & \textbf{28.4} & \textbf{0.912} & \textbf{14.6} & \textbf{31.2} \\
\hline
\end{tabular}
\end{table}

Each component contributes incrementally, with recognition feature loss providing the largest improvement in HTR metrics while maintaining visual quality.

\subsubsection{Frozen vs. Joint Recognizer Training}

% [PLACEHOLDER: Table 8 - Frozen vs. joint recognizer comparison]
\begin{table}[!t]
\renewcommand{\arraystretch}{1.3}
\caption{Frozen vs. Joint Recognizer Training Comparison}
\label{table_ablation_recognizer}
\centering
\begin{tabular}{|l|c|c|c|}
\hline
\textbf{Recognizer Training} & \textbf{SSIM} & \textbf{CER (\%)} & \textbf{Stability} \\
\hline
Joint (trainable) & 0.878 & 19.3 & Unstable \\
Frozen (ours) & \textbf{0.912} & \textbf{14.6} & Stable \\
\hline
\end{tabular}
\end{table}

Frozen recognizer approach achieves better performance and training stability, supporting our design choice.

\subsection{Qualitative Analysis}

Figure~\ref{fig:qualitative_results} presents visual comparisons on challenging examples.

% [PLACEHOLDER: Figure 6 - Qualitative visual comparison]
\begin{figure*}[!t]
\centering
% \includegraphics[width=7in]{fig_qualitative_results}
\fbox{\parbox{7in}{\centering [PLACEHOLDER FIGURE 6]\\Qualitative Results Comparison\\Rows: (a) Degraded input, (b) Sauvola, (c) Standard GAN\\(d) DE-GAN, (e) HTR-GAN, (f) Ours, (g) Ground Truth\\Multiple examples showing text preservation and noise removal}}
\caption{Qualitative comparison of restoration results on challenging degraded samples. Our method better preserves thin strokes and character connectivity while removing background noise.}
\label{fig:qualitative_results}
\end{figure*}

\textbf{Observations:}
\begin{itemize}
    \item Classical methods (Sauvola) produce harsh binarization, losing thin strokes
    \item Standard GAN removes noise effectively but sometimes breaks character connections
    \item DE-GAN produces clean results but may over-smooth fine details
    \item HTR-GAN preserves text better but shows some remaining noise
    \item Our method achieves optimal balance: clean background with intact, coherent text
\end{itemize}

\subsection{Failure Case Analysis}

We analyze failure cases to identify limitations:

\begin{enumerate}
    \item \textbf{Extremely faded text:} When original text is nearly invisible, restoration may hallucinate strokes
    \item \textbf{Overlapping text/stamps:} Separating overlapping content remains challenging
    \item \textbf{Out-of-vocabulary scripts:} Paleographic abbreviations or symbols not seen in training
\end{enumerate}

Figure~\ref{fig:failure_cases} shows representative failure examples.

% [PLACEHOLDER: Figure 7 - Failure case examples]
\begin{figure}[!t]
\centering
% \includegraphics[width=3.5in]{fig_failure_cases}
\fbox{\parbox{3.5in}{\centering [PLACEHOLDER FIGURE 7]\\Failure Cases\\(a) Extremely faded text - hallucination\\(b) Overlapping stamp - incomplete separation\\(c) Paleographic symbols - misreconstruction}}
\caption{Representative failure cases and limitations of the proposed method.}
\label{fig:failure_cases}
\end{figure}

\subsection{Computational Efficiency}

Table~\ref{table_efficiency} compares computational requirements.

% [PLACEHOLDER: Table 9 - Computational efficiency comparison]
\begin{table}[!t]
\renewcommand{\arraystretch}{1.3}
\caption{Computational Efficiency Comparison}
\label{table_efficiency}
\centering
\begin{tabular}{|l|c|c|c|}
\hline
\textbf{Method} & \textbf{Params (M)} & \textbf{Inference (ms)} & \textbf{GPU Memory (GB)} \\
\hline
Sauvola & N/A & 8 & 0.1 \\
U-Net & 31.2 & 18 & 2.3 \\
Standard GAN & 34.7 & 22 & 2.8 \\
DE-GAN & 38.1 & 25 & 3.1 \\
HTR-GAN & 45.3 & 35 & 4.2 \\
\textbf{Ours} & 47.8 & 37 & 4.5 \\
\hline
\end{tabular}
\end{table}

Our method has slightly higher computational cost due to dual-modal discriminator and frozen recognizer, but remains practical for batch processing (37ms per image on GPU).

% ============================================================
% VI. DISCUSSION
% ============================================================
\section{Discussion}

This section discusses key findings, design choices, practical implications, and future directions.

\subsection{Impact of Dual-Modal Discriminator}

The dual-modal discriminator represents a key innovation. Our ablation studies (Table~\ref{table_ablation_discriminator}) demonstrate that combining CNN and LSTM pathways yields superior results compared to single-pathway variants. The LSTM pathway explicitly models sequential text structure, enabling detection of subtle text coherence issues that purely spatial discriminators miss.

Qualitative analysis reveals that the dual-modal discriminator helps maintain character connectivity and stroke continuity, particularly for cursive handwriting where characters flow together. This is evident in examples where single-pathway discriminators break ligatures while our method preserves them.

\subsection{Frozen vs. Joint Recognizer Training}

Our choice to use a frozen pre-trained recognizer rather than joint training (as in~\cite{souibgui2021enhance}) proves beneficial. Joint training introduces optimization challenges: the recognizer and generator can evolve in conflicting directions, and recognizer degradation during training provides unstable gradients.

The frozen approach provides consistent, reliable text-aware feedback throughout training. By pre-training the recognizer on diverse clean handwriting datasets, we encode robust text feature representations that guide restoration without requiring text labels during GAN training.

\subsection{Loss Weight Optimization Insights}

Our Bayesian optimization and fANOVA analysis (Table~\ref{table_param_importance}) reveal that adversarial loss weight $\lambda_{\text{adv}}$ is the most critical hyperparameter (66.4\% importance). This aligns with GAN theory: adversarial loss governs the realism-fidelity tradeoff and training dynamics.

Too high $\lambda_{\text{adv}}$ causes mode collapse (generator produces limited variety) and training instability. Too low $\lambda_{\text{adv}}$ results in blurry outputs lacking realism. Our optimized value (2.5) balances these concerns.

Recognition feature loss weight $\lambda_{\text{rec}}$ (24.3\% importance) has moderate impact, suggesting text-aware guidance is valuable but must be balanced with visual quality objectives.

\subsection{Generalization to Real Historical Documents}

Performance on real ANRI documents (Table~\ref{table_anri_results}) demonstrates reasonable generalization despite training only on synthetic degradation. However, the performance gap (CER: 14.6\% synthetic vs. 21.3\% real) indicates room for improvement.

Potential strategies for better generalization:
\begin{enumerate}
    \item \textbf{Domain adaptation:} Fine-tune on small sets of real degraded documents
    \item \textbf{Improved synthetic degradation:} Incorporate more realistic aging patterns
    \item \textbf{Unsupervised/self-supervised learning:} Leverage unlabeled real documents
\end{enumerate}

\subsection{Practical Deployment Considerations}

For deploying our system in archival digitization workflows, several considerations arise:

\textbf{Batch Processing:} Our method processes images at 37ms/image (GPU), enabling throughput of $\sim$27 images/second, suitable for large-scale archival projects.

\textbf{Quality Control:} Despite strong average performance, failure cases exist (Figure~\ref{fig:failure_cases}). We recommend human review for critical documents or implementing confidence scoring to flag uncertain restorations.

\textbf{Preservation Ethics:} Document restoration modifies original appearance. For scholarly use, both original and restored versions should be archived, with clear provenance metadata.

\subsection{Limitations and Future Work}

\textbf{Current Limitations:}
\begin{enumerate}
    \item Requires line-level segmentation as preprocessing
    \item Struggles with extremely faded or damaged text
    \item Limited to grayscale; color documents require adaptation
    \item Trained primarily on Latin/Dutch scripts; multilingual generalization untested
\end{enumerate}

\textbf{Future Research Directions:}
\begin{enumerate}
    \item \textbf{End-to-end page-level restoration:} Extend to full page images without line segmentation
    \item \textbf{Multi-task learning:} Joint restoration and transcription in unified framework
    \item \textbf{Uncertainty quantification:} Provide confidence scores for restoration quality
    \item \textbf{Interactive restoration:} Allow user feedback to guide difficult regions
    \item \textbf{Cross-lingual evaluation:} Test on Arabic, Chinese, and other scripts
    \item \textbf{Few-shot adaptation:} Enable rapid adaptation to new document types with minimal examples
\end{enumerate}

% ============================================================
% VII. CONCLUSION
% ============================================================
\section{Conclusion}

This paper presented a novel HTR-oriented document restoration framework addressing a critical gap in existing document enhancement methods: the lack of explicit optimization for text recognition performance. Our approach integrates three key innovations within a GAN architecture:

\begin{enumerate}
    \item A dual-modal discriminator with parallel CNN and LSTM pathways that simultaneously assesses visual quality and sequential text coherence
    \item A frozen pre-trained HTR recognizer providing stable text-aware gradients during training without requiring text labels
    \item A systematically optimized multi-component loss function balancing adversarial, reconstruction, perceptual, and recognition objectives
\end{enumerate}

Extensive experiments on synthetic degraded datasets and real 16th-18th century historical manuscripts from the Indonesian National Archives demonstrate significant improvements over state-of-the-art methods. Our approach achieves PSNR of 28.42 dB, SSIM of 0.912, and reduces Character Error Rate to 14.6\% on synthetic test data - a 24.3\% relative reduction compared to the best baseline. On real historical documents, we achieve 21.3\% CER, substantially outperforming all comparison methods.

Ablation studies confirm that each architectural component contributes to overall performance, with the dual-modal discriminator and recognition feature loss providing the most substantial improvements in HTR metrics. Our analysis reveals that adversarial loss weight is the most critical hyperparameter, requiring careful calibration through systematic optimization.

The proposed framework demonstrates that explicitly incorporating HTR objectives during document restoration training leads to outputs that are both visually cleaner and more readable by automated transcription systems. This has important implications for large-scale historical document digitization projects, where the ultimate goal is machine-readable text extraction rather than solely visual quality improvement.

While our method shows strong performance, several limitations remain, including dependence on line segmentation, challenges with extremely degraded content, and the need for better generalization to diverse scripts and degradation types. Future work will address these limitations and extend the framework to page-level processing and multi-lingual scenarios.

In conclusion, this work advances the state-of-the-art in document restoration by bridging the gap between visual enhancement and text recognition, providing a practical and effective solution for preserving and making accessible the world's invaluable historical document collections.

% ============================================================
% APPENDICES (if needed)
% ============================================================
\appendices

\section{Network Architecture Details}

% [PLACEHOLDER: Detailed layer-by-layer architecture specifications]

This appendix provides detailed layer-by-layer specifications for reproducibility.

\subsection{Generator Architecture}

Detailed U-Net architecture:
\begin{itemize}
    \item Input: $64 \times 512 \times 1$ (grayscale line image)
    \item Encoder blocks: [Conv-BN-ReLU-Conv-BN-ReLU-MaxPool]
    \item Channel progression: 64, 128, 256, 512, 512, 512, 512, 512
    \item Bottleneck: 512 channels
    \item Decoder blocks: [UpConv-Concat-Conv-BN-ReLU-Conv-BN-ReLU]
    \item Output: $64 \times 512 \times 1$ with sigmoid activation
\end{itemize}

\subsection{Dual-Modal Discriminator Architecture}

CNN pathway specifications:
\begin{itemize}
    \item Input: Concatenated degraded + generated/real ($64 \times 512 \times 2$)
    \item Layer 1: Conv(64, 4×4, stride=2) - LeakyReLU(0.2)
    \item Layer 2: Conv(128, 4×4, stride=2) - BN - LeakyReLU(0.2)
    \item Layer 3: Conv(256, 4×4, stride=2) - BN - LeakyReLU(0.2)
    \item Layer 4: Conv(512, 4×4, stride=2) - BN - LeakyReLU(0.2)
\end{itemize}

LSTM pathway specifications:
\begin{itemize}
    \item Reshape CNN features to sequence format
    \item Bidirectional LSTM: 256 units forward + 256 backward
    \item Output: 512-dimensional sequential features
\end{itemize}

Fusion and classification:
\begin{itemize}
    \item Concatenate CNN and LSTM outputs
    \item FC layer: 512 units - LeakyReLU(0.2) - Dropout(0.3)
    \item FC layer: 256 units - LeakyReLU(0.2)
    \item Output layer: 1 unit - Sigmoid
\end{itemize}

\section{Training Hyperparameters}

% [PLACEHOLDER: Complete list of hyperparameters]

\begin{table}[!h]
\renewcommand{\arraystretch}{1.3}
\caption{Complete Hyperparameter Configuration}
\label{table_hyperparams}
\centering
\begin{tabular}{|l|c|}
\hline
\textbf{Hyperparameter} & \textbf{Value} \\
\hline
Learning rate (G) & $2 \times 10^{-4}$ \\
Learning rate (D) & $2 \times 10^{-4}$ \\
Optimizer & Adam \\
$\beta_1$ & 0.5 \\
$\beta_2$ & 0.999 \\
Batch size & 16 \\
Epochs (main training) & 100 \\
Epochs (fine-tuning) & 50 \\
Image size (H × W) & $64 \times 512$ \\
Dropout rate & 0.3 \\
Label smoothing range & [0.9, 1.0] \\
Gradient clipping norm & 1.0 \\
$\lambda_{\text{adv}}$ & 2.5 \\
$\lambda_{\text{pixel}}$ & 120.0 \\
$\lambda_{\text{perc}}$ & 10.0 \\
$\lambda_{\text{rec}}$ & 80.0 \\
\hline
\end{tabular}
\end{table}

\section{Additional Experimental Results}

% [PLACEHOLDER: Additional tables, graphs, etc.]

\subsection{Per-Degradation-Level Analysis}

We analyze performance across different degradation severity levels:

% [PLACEHOLDER: Table showing performance vs. degradation level]

\subsection{Statistical Significance Tests}

We perform paired t-tests comparing our method with baselines. All improvements are statistically significant with $p < 0.001$.

% ============================================================
% ACKNOWLEDGMENT
% ============================================================
\section*{Acknowledgment}

The authors would like to thank the Indonesian National Archives (Arsip Nasional Republik Indonesia) for providing access to historical document collections and paleography expertise. We also acknowledge [University/Institution] for computational resources and [Funding Agency] for financial support under grant [Grant Number].

% ============================================================
% REFERENCES
% ============================================================
\begin{thebibliography}{99}

\bibitem{otsu1979threshold}
N. Otsu, ``A threshold selection method from gray-level histograms,'' \textit{IEEE Trans. Syst., Man, Cybern.}, vol. 9, no. 1, pp. 62--66, Jan. 1979.

\bibitem{sauvola2000adaptive}
J. Sauvola and M. Pietikäinen, ``Adaptive document image binarization,'' \textit{Pattern Recognit.}, vol. 33, no. 2, pp. 225--236, Feb. 2000.

\bibitem{niblack1986introduction}
W. Niblack, \textit{An Introduction to Digital Image Processing}. Englewood Cliffs, NJ, USA: Prentice-Hall, 1986.

\bibitem{wolf2004extraction}
C. Wolf, J.-M. Jolion, and F. Chassaing, ``Text localization, enhancement and binarization in multimedia documents,'' in \textit{Proc. 16th Int. Conf. Pattern Recognit.}, vol. 2, 2002, pp. 1037--1040.

\bibitem{tensmeyer2017document}
C. Tensmeyer and T. Martinez, ``Document image binarization with fully convolutional neural networks,'' in \textit{Proc. 14th IAPR Int. Conf. Document Anal. Recognit.}, vol. 1, 2017, pp. 99--104.

\bibitem{calvo2018selectional}
A. Calvo-Zaragoza and J. Oncina, ``Recognition of pen-based music notation: The HOMUS dataset,'' in \textit{Proc. 22nd Int. Conf. Pattern Recognit.}, 2014, pp. 3038--3043.

\bibitem{zhao2019document}
J. Zhao, C. Shi, F. Jia, Y. Wang, and B. Xiao, ``Document image binarization with cascaded generators of conditional generative adversarial networks,'' \textit{Pattern Recognit.}, vol. 96, p. 106968, Dec. 2019.

\bibitem{vo2018binarization}
Q. N. Vo, S. H. Kim, H. J. Yang, and G. Lee, ``Binarization of degraded document images based on hierarchical deep supervised network,'' \textit{Pattern Recognit.}, vol. 74, pp. 568--586, Feb. 2018.

\bibitem{souibgui2020degan}
M. A. Souibgui, S. Khamekhem Jemni, Y. Kessentini, C. Fornés, A. Fornés, and C. Ogier, ``DE-GAN: A conditional generative adversarial network for document enhancement,'' \textit{IEEE Trans. Pattern Anal. Mach. Intell.}, vol. 44, no. 3, pp. 1180--1191, Mar. 2022.

\bibitem{souibgui2021enhance}
M. A. Souibgui, S. Khamekhem Jemni, Y. Kessentini, and A. Fornés, ``Enhance to read better: A multi-task adversarial network for handwritten document image enhancement,'' \textit{Pattern Recognit.}, vol. 123, p. 108368, Mar. 2022.

\bibitem{ronneberger2015unet}
O. Ronneberger, P. Fischer, and T. Brox, ``U-Net: Convolutional networks for biomedical image segmentation,'' in \textit{Proc. Int. Conf. Med. Image Comput. Comput.-Assist. Interv.}, 2015, pp. 234--241.

\bibitem{johnson2016perceptual}
J. Johnson, A. Alahi, and L. Fei-Fei, ``Perceptual losses for real-time style transfer and super-resolution,'' in \textit{Proc. Eur. Conf. Comput. Vis.}, 2016, pp. 694--711.

\bibitem{wang2004ssim}
Z. Wang, A. C. Bovik, H. R. Sheikh, and E. P. Simoncelli, ``Image quality assessment: From error visibility to structural similarity,'' \textit{IEEE Trans. Image Process.}, vol. 13, no. 4, pp. 600--612, Apr. 2004.

\bibitem{li2021trocr}
M. Li et al., ``TrOCR: Transformer-based optical character recognition with pre-trained models,'' in \textit{Proc. AAAI Conf. Artif. Intell.}, vol. 37, 2023, pp. 13094--13102.

\bibitem{marti2002iam}
U.-V. Marti and H. Bunke, ``The IAM-database: An English sentence database for offline handwriting recognition,'' \textit{Int. J. Document Anal. Recognit.}, vol. 5, no. 1, pp. 39--46, Nov. 2002.

\bibitem{bergstra2011algorithms}
J. Bergstra, R. Bardenet, Y. Bengio, and B. Kégl, ``Algorithms for hyper-parameter optimization,'' in \textit{Proc. Adv. Neural Inf. Process. Syst.}, 2011, pp. 2546--2554.

\bibitem{akiba2019optuna}
T. Akiba, S. Sano, T. Yanase, T. Ohta, and M. Koyama, ``Optuna: A next-generation hyperparameter optimization framework,'' in \textit{Proc. 25th ACM SIGKDD Int. Conf. Knowl. Discovery Data Mining}, 2019, pp. 2623--2631.

% [ADD MORE REFERENCES AS NEEDED]

\end{thebibliography}

% ============================================================
% BIOGRAPHIES
% ============================================================

\begin{IEEEbiographynophoto}{[Your Name]}
received the B.S. degree in [Your Field] from [Your University] in [Year]. Currently pursuing Ph.D. degree in [Your Department] at [Your University]. Research interests include document image analysis, generative adversarial networks, and historical document preservation.
\end{IEEEbiographynophoto}

\begin{IEEEbiographynophoto}{[Supervisor Name]}
is a Professor at [Department], [University]. Research interests include computer vision, pattern recognition, and machine learning. Published over [Number] papers in top-tier conferences and journals.
\end{IEEEbiographynophoto}

\begin{IEEEbiographynophoto}{[Co-Supervisor Name]}
is an Associate Professor at [Department], [University]. Specializes in document analysis, OCR/HTR systems, and deep learning applications for cultural heritage preservation.
\end{IEEEbiographynophoto}

\end{document}


% An example of a floating figure using the graphicx package.
% Note that \label must occur AFTER (or within) \caption.
% For figures, \caption should occur after the \includegraphics.
% Note that IEEEtran v1.7 and later has special internal code that
% is designed to preserve the operation of \label within \caption
% even when the captionsoff option is in effect. However, because
% of issues like this, it may be the safest practice to put all your
% \label just after \caption rather than within \caption{}.
%
% Reminder: the "draftcls" or "draftclsnofoot", not "draft", class
% option should be used if it is desired that the figures are to be
% displayed while in draft mode.
%
%\begin{figure}[!t]
%\centering
%\includegraphics[width=2.5in]{myfigure}
% where an .eps filename suffix will be assumed under latex, 
% and a .pdf suffix will be assumed for pdflatex; or what has been declared
% via \DeclareGraphicsExtensions.
%\caption{Simulation results for the network.}
%\label{fig_sim}
%\end{figure}

% Note that the IEEE typically puts floats only at the top, even when this
% results in a large percentage of a column being occupied by floats.


% An example of a double column floating figure using two subfigures.
% (The subfig.sty package must be loaded for this to work.)
% The subfigure \label commands are set within each subfloat command,
% and the \label for the overall figure must come after \caption.
% \hfil is used as a separator to get equal spacing.
% Watch out that the combined width of all the subfigures on a 
% line do not exceed the text width or a line break will occur.
%
%\begin{figure*}[!t]
%\centering
%\subfloat[Case I]{\includegraphics[width=2.5in]{box}%
%\label{fig_first_case}}
%\hfil
%\subfloat[Case II]{\includegraphics[width=2.5in]{box}%
%\label{fig_second_case}}
%\caption{Simulation results for the network.}
%\label{fig_sim}
%\end{figure*}
%
% Note that often IEEE papers with subfigures do not employ subfigure
% captions (using the optional argument to \subfloat[]), but instead will
% reference/describe all of them (a), (b), etc., within the main caption.
% Be aware that for subfig.sty to generate the (a), (b), etc., subfigure
% labels, the optional argument to \subfloat must be present. If a
% subcaption is not desired, just leave its contents blank,
% e.g., \subfloat[].


% An example of a floating table. Note that, for IEEE style tables, the
% \caption command should come BEFORE the table and, given that table
% captions serve much like titles, are usually capitalized except for words
% such as a, an, and, as, at, but, by, for, in, nor, of, on, or, the, to
% and up, which are usually not capitalized unless they are the first or
% last word of the caption. Table text will default to \footnotesize as
% the IEEE normally uses this smaller font for tables.
% The \label must come after \caption as always.
%
%\begin{table}[!t]
%% increase table row spacing, adjust to taste
%\renewcommand{\arraystretch}{1.3}
% if using array.sty, it might be a good idea to tweak the value of
% \extrarowheight as needed to properly center the text within the cells
%\caption{An Example of a Table}
%\label{table_example}
%\centering
%% Some packages, such as MDW tools, offer better commands for making tables
%% than the plain LaTeX2e tabular which is used here.
%\begin{tabular}{|c||c|}
%\hline
%One & Two\\
%\hline
%Three & Four\\
%\hline
%\end{tabular}
%\end{table}


% Note that the IEEE does not put floats in the very first column
% - or typically anywhere on the first page for that matter. Also,
% in-text middle ("here") positioning is typically not used, but it
% is allowed and encouraged for Computer Society conferences (but
% not Computer Society journals). Most IEEE journals/conferences use
% top floats exclusively. 
% Note that, LaTeX2e, unlike IEEE journals/conferences, places
% footnotes above bottom floats. This can be corrected via the
% \fnbelowfloat command of the stfloats package.




\section{Conclusion}
The conclusion goes here.





% if have a single appendix:
%\appendix[Proof of the Zonklar Equations]
% or
%\appendix  % for no appendix heading
% do not use \section anymore after \appendix, only \section*
% is possibly needed

% use appendices with more than one appendix
% then use \section to start each appendix
% you must declare a \section before using any
% \subsection or using \label (\appendices by itself
% starts a section numbered zero.)
%


\appendices
\section{Proof of the First Zonklar Equation}
Appendix one text goes here.

% you can choose not to have a title for an appendix
% if you want by leaving the argument blank
\section{}
Appendix two text goes here.


% use section* for acknowledgment
\section*{Acknowledgment}


The authors would like to thank...


% Can use something like this to put references on a page
% by themselves when using endfloat and the captionsoff option.
\ifCLASSOPTIONcaptionsoff
  \newpage
\fi



% trigger a \newpage just before the given reference
% number - used to balance the columns on the last page
% adjust value as needed - may need to be readjusted if
% the document is modified later
%\IEEEtriggeratref{8}
% The "triggered" command can be changed if desired:
%\IEEEtriggercmd{\enlargethispage{-5in}}

% references section

% can use a bibliography generated by BibTeX as a .bbl file
% BibTeX documentation can be easily obtained at:
% http://mirror.ctan.org/biblio/bibtex/contrib/doc/
% The IEEEtran BibTeX style support page is at:
% http://www.michaelshell.org/tex/ieeetran/bibtex/
%\bibliographystyle{IEEEtran}
% argument is your BibTeX string definitions and bibliography database(s)
%\bibliography{IEEEabrv,../bib/paper}
%
% <OR> manually copy in the resultant .bbl file
% set second argument of \begin to the number of references
% (used to reserve space for the reference number labels box)
\begin{thebibliography}{1}

\bibitem{IEEEhowto:kopka}
H.~Kopka and P.~W. Daly, \emph{A Guide to \LaTeX}, 3rd~ed.\hskip 1em plus
  0.5em minus 0.4em\relax Harlow, England: Addison-Wesley, 1999.

\end{thebibliography}

% biography section
% 
% If you have an EPS/PDF photo (graphicx package needed) extra braces are
% needed around the contents of the optional argument to biography to prevent
% the LaTeX parser from getting confused when it sees the complicated
% \includegraphics command within an optional argument. (You could create
% your own custom macro containing the \includegraphics command to make things
% simpler here.)
%\begin{IEEEbiography}[{\includegraphics[width=1in,height=1.25in,clip,keepaspectratio]{mshell}}]{Michael Shell}
% or if you just want to reserve a space for a photo:

\begin{IEEEbiography}{Michael Shell}
Biography text here.
\end{IEEEbiography}

% if you will not have a photo at all:
\begin{IEEEbiographynophoto}{John Doe}
Biography text here.
\end{IEEEbiographynophoto}

% insert where needed to balance the two columns on the last page with
% biographies
%\newpage

\begin{IEEEbiographynophoto}{Jane Doe}
Biography text here.
\end{IEEEbiographynophoto}

% You can push biographies down or up by placing
% a \vfill before or after them. The appropriate
% use of \vfill depends on what kind of text is
% on the last page and whether or not the columns
% are being equalized.

%\vfill

% Can be used to pull up biographies so that the bottom of the last one
% is flush with the other column.
%\enlargethispage{-5in}



% that's all folks
\end{document}